\section{Introduction}
\label{introduction}
%%% intro problema fisico: perché è difficile predirre il comportamento delle scie eoliche? perché è importante riuscire a farlo? spiega natura del fenomeno e difficoltà di modellizzazione
%%% Spiega che esistono diverse tecniche per modellizzare le scie eoliche, ma che tutte hanno dei limiti (o fedeltà o costo computazionale).
%%% cita l uso di tecniche di machine learning in campo eolico per la predizione della potenza, gestione manutentiva, ottimizzazione del layout
%%% spiega che in questo lavoro ci si concentra su un approccio di machine learning per la predizione del comportamento delle scie eoliche, quindi un qualcosa di diverso rispetto a quanto visto perché "si tratta proprio di creare da zero un modello di turbolenza ML-based" e non di usare un modello di turbolenza esistente per predire la potenza o ottimizzare il layout e poi usare tecniche di ML per migliorare la predizione o usarle in post sulla base di quei risultati per ottimizzare il layout o la gestione manutentiva.
%%% fai esempi di lavori che hanno usato tecniche di ML per predire la potenza o ottimizzare il layout e anche di lavori che hanno usato tecniche di ML per migliorare la predizione dei modelli di turbolenza esistenti. 
%%% spiega che questo approccio è interessante perché potrebbe permettere di superare i limiti dei modelli di turbolenza tradizionali, che spesso non riescono a catturare la complessità del fenomeno eolico, e potrebbe portare a predizioni più accurate e a una migliore comprensione del comportamento delle scie eoliche nonché a una possibile futura applicazione in real time con il funzionamento della farm attraveso la possibilità di ottenere modelli rappresentiativi compatti del campo e del suo funzionamento. 
Wind energy is one of the cornerstones of the ongoing transition toward sustainable electricity generation. As wind turbines are deployed in increasingly large farms, both onshore and offshore, the aerodynamic interaction among turbines becomes a critical design and operational challenge. When a turbine extracts kinetic energy from the incoming flow, it leaves behind a wake region characterised by a significant velocity deficit and elevated turbulence levels. Downstream turbines that operate within this region consequently suffer both reduced power output and increased fatigue loads, affecting the economic and structural performance of the entire farm~\cite{Porte-Agel2020}. Accurate prediction of wind turbine wake behaviour is therefore essential for wind farm layout optimisation, energy yield estimation, and structural load assessment.

From a fluid-mechanics perspective, the wake of a horizontal-axis wind turbine (HAWT) is a complex, multi-scale phenomenon. It is deeply coupled with the Atmospheric Boundary Layer (ABL), which is itself inherently unsteady and stratified~\cite{Porte-Agel2020}. The near-wake region, dominated by the tip vortex system shed by the rotor blades, transitions into a far-wake where shear-driven turbulent mixing progressively restores the velocity deficit. This recovery process is strongly influenced by the ambient turbulence intensity, wind shear, thermal stability, and the topography of the terrain~\cite{heck2025unravelingeffectsatmosphericdynamics}. The wide range of length and time scales involved, from blade boundary layers to the full extent of the ABL, makes any single simulation approach inherently challenging.

Several modelling strategies are available, each representing a different trade-off between physical fidelity and computational cost. On one end of the spectrum, analytical engineering wake models (e.g., the Jensen top-hat model and its Gaussian successors) offer near-instantaneous evaluation and are therefore routinely used in wind farm layout optimisation~\cite{GholamiAnjiraki2025}, but they rely on strong simplifying assumptions and cannot capture the complex turbulent structure of real wakes. On the other end, Large Eddy Simulation (LES) resolves the dominant turbulent structures explicitly and is considered the reference standard for wake physics; however, a single simulation of a medium-sized wind farm can demand between $10^5$ and $3 \times 10^6$ CPU hours~\cite{Steiner2021}, making it prohibitively expensive for systematic design or optimisation studies. Unsteady Reynolds-Averaged Navier-Stokes (URANS) methods occupy an intermediate position, reducing computational cost by roughly two orders of magnitude relative to LES~\cite{Steiner2021}. Yet the closure assumptions embedded in standard RANS turbulence models, most notably the Boussinesq eddy-viscosity hypothesis and the isotropic treatment of the Reynolds stress tensor, introduce structural model-form errors that are particularly pronounced in separated and wake flows, leading to a systematic mis-prediction of the wake recovery rate~\cite{72116acdc5aa4b3b8fc752f1b2fa9958,Moss2023}.

Machine Learning (ML) has emerged as a powerful paradigm to complement or enhance physics-based simulations across several domains of wind energy research. A large body of literature has applied ML to operational tasks such as wind power forecasting from SCADA data~\cite{Infield2022,Moss2023}, predictive maintenance and condition monitoring of turbine components~\cite{Infield2022}, and wind farm layout optimisation using data-driven surrogate models of wake interactions~\cite{GholamiAnjiraki2025,YANG2023119240}. In these applications ML is generally used downstream of a physical model: a CFD or analytical solver first generates the flow field, and ML then post-processes or surrogate-replaces its output to predict power, loads, or optimal configurations. 

A conceptually distinct and more recent line of research aims instead to use ML to correct or replace the turbulence closure itself, modifying the governing equations during the simulation. An example of that can be seen in ~\cite{Duraisamy2019} with the proposed FIML (field inversion and machine learning) proposed by Parish \& Duraisamy and the invariant neural-network closure of Ling et al.~\cite{Ling2016}. In the wind energy context, Steiner et al.~\cite{Steiner2021} demonstrated that optimal corrective fields for the turbulence anisotropy tensor and turbulent kinetic energy production can be derived from time-averaged LES data and used to construct a custom RANS closure via sparse symbolic regression, yielding substantially improved wake predictions for multi-turbine configurations at wind-tunnel scale. Similarly, data-driven Gaussian-process corrections to the eddy-viscosity field, trained on single-turbine LES, have shown improved accuracy when deployed in a wind-plant RANS simulation~\cite{King_2018}. More broadly, Steiner~\cite{72116acdc5aa4b3b8fc752f1b2fa9958} showed that combined baseline-RANS plus data-driven correction strategies can outperform standard models for both velocity and turbulent kinetic energy, while acknowledging that larger datasets and more stable numerical implementations are needed before such methods can be applied at industrial scale.

The present work contributes to this emerging research direction, with specific novelty in three respects. First, rather than identifying a global correction to the turbulence model, a clustering algorithm (CA) based on the K-means method is employed to segment the flow domain into physically distinct regions characterised by different wake dynamics, following the decomposition approach introduced in~\cite{Clustering_tieghi,Carloni2026}. This region-aware strategy naturally improves the generalisability and interpretability of the trained models and makes the correction less sensitive to mesh refinement. Second, a Multi-Layer Perceptron (MLP) is trained to predict the local correction $\Delta\nu_t = \nu_{t,eq} - \nu_{t,\mathrm{RANS}}$ between the LES-equivalent eddy viscosity and the RANS baseline, and the resulting data-driven closure is directly embedded into the OpenFOAM PIMPLE solver via a modified turbulence correction loop, enabling its use in fully predictive RANS simulations. Third, the methodology is tested across multiple mesh refinement levels, demonstrating that LES-quality accuracy in wake recovery can be achieved at a computational cost lower than that of standard LES, which is a prerequisite for practical wind farm applications.

It is worth emphasising what distinguishes this approach from the bulk of ML applications in wind energy. Most existing methods employ ML as a post-processing or surrogate tool applied to the output of a turbulence model; the turbulence model itself is left unchanged. Here, the ML component is used to define the turbulence model, to construct, from scratch, a problem-specific eddy-viscosity closure calibrated to reproduce high-fidelity LES statistics. This opens a path towards compact, computationally efficient flow models that could in the future support real-time wind farm monitoring and control, by providing accurate yet computational effective representations of the wake field and its evolution under varying inflow conditions.

The paper is structured as follows. Sec.~\ref{sec:methodology} describes the CFD campaign used to generate the training and testing data, the clustering pipeline, and the MLP architecture and training strategy. Sec.~\ref{sec:case_study} introduces the NREL-5MW reference turbine case study and the computational domain. Results are presented in Sec.~\ref{sec:results}, covering the clustering, the neural network training metrics, and the performance of the ML-based RANS model on coarser grids. Conclusions are drawn in Sec.~\ref{sec:conclusions}.
